% wave13_a13_hCS_to_delta5_costello.tex
% Wave 13 / A13 --- Costello channel, adversarial pass
% Mission: 6d hCS on K3 x E --> E_3 holomorphic factorisation algebra -->
% Sp_{K3, E} (integrate out K3, restrict to E) --> Lie-envelope chiral avatar;
% compare against Lorgat 2020 explicit construction of g_{Delta_5} and Delta_5.
% Cross-refs: wave12_a13_costello_5_routes_6d_hCS.tex; wave12_b3_6d_hCS_E3_gen_rel.tex;
% wave12_b4_6d_hCS_BV_BRST_Feynman.tex; wave12_u1_two_stage_functor.tex.
% Source paper: notes/wave13_source_paper.txt = Lorgat 2020 (arXiv: user's own).

\documentclass[11pt]{amsart}
\usepackage[T1]{fontenc}
\usepackage[utf8]{inputenc}
\usepackage{amsmath,amssymb,amsthm}
\usepackage{mathtools}
\usepackage[margin=1.1in]{geometry}
\usepackage{hyperref}

\providecommand{\C}{\mathbb{C}}
\providecommand{\R}{\mathbb{R}}
\providecommand{\Z}{\mathbb{Z}}
\providecommand{\N}{\mathbb{N}}
\providecommand{\Q}{\mathbb{Q}}
\providecommand{\cE}{\mathcal{E}}
\providecommand{\cF}{\mathcal{F}}
\providecommand{\cO}{\mathcal{O}}
\providecommand{\cA}{\mathcal{A}}
\providecommand{\cZ}{\mathcal{Z}}
\providecommand{\cC}{\mathcal{C}}
\providecommand{\cH}{\mathcal{H}}
\providecommand{\cL}{\mathcal{L}}
\providecommand{\cW}{\mathcal{W}}
\providecommand{\dbar}{\bar{\partial}}
\providecommand{\g}{\mathfrak{g}}
\providecommand{\h}{\mathfrak{h}}
\providecommand{\fsl}{\mathfrak{sl}}
\providecommand{\fu}{\mathfrak{u}}
\providecommand{\tr}{\mathrm{tr}}
\providecommand{\Tr}{\mathrm{Tr}}
\providecommand{\Obs}{\mathrm{Obs}}
\providecommand{\Sym}{\mathrm{Sym}}
\providecommand{\End}{\mathrm{End}}
\providecommand{\Hom}{\mathrm{Hom}}
\providecommand{\id}{\mathrm{id}}
\providecommand{\hCS}{\mathrm{hCS}}
\providecommand{\BV}{\mathrm{BV}}
\providecommand{\Fact}{\mathrm{Fact}}
\providecommand{\PreFact}{\mathrm{PreFact}}
\providecommand{\E}{E}
\providecommand{\Eone}{E_{1}}
\providecommand{\Etwo}{E_{2}}
\providecommand{\Ethree}{E_{3}}
\providecommand{\Disk}{\mathrm{Disk}}
\providecommand{\Ran}{\mathrm{Ran}}
\providecommand{\CE}{\mathrm{CE}}
\providecommand{\Heis}{\mathsf{Heis}}
\providecommand{\Muk}{\Lambda_{\mathrm{Muk}}}
\providecommand{\HB}{\Lambda^{(2,1)}_{I\!I}}
\providecommand{\HA}{\Lambda^{(2,1)}}
\providecommand{\LamF}{\Lambda^{(3,2)}}
\providecommand{\Theta}{\mathbf{\Theta}}
\providecommand{\Gdelta}{\mathfrak{g}_{\Delta_{5}}}
\providecommand{\PhiFA}{\Phi^{\mathrm{FA}}}
\providecommand{\SpF}{\mathrm{Sp}}
\providecommand{\kch}{\kappa_{\mathrm{ch}}}
\providecommand{\kcat}{\kappa_{\mathrm{cat}}}
\providecommand{\kBKM}{\kappa_{\mathrm{BKM}}}
\providecommand{\kfib}{\kappa_{\mathrm{fiber}}}
\providecommand{\ClaimStatusTheorem}{\textsc{[theorem]}}
\providecommand{\ClaimStatusConjectured}{\textsc{[conjectural]}}
\providecommand{\ClaimStatusHeuristic}{\textsc{[heuristic]}}
\providecommand{\ClaimStatusDerivation}{\textsc{[derivation]}}
\providecommand{\ClaimStatusDefinition}{\textsc{[definition]}}

\theoremstyle{plain}
\newtheorem{theorem}{Theorem}[section]
\newtheorem{proposition}[theorem]{Proposition}
\newtheorem{lemma}[theorem]{Lemma}
\newtheorem{corollary}[theorem]{Corollary}
\newtheorem{conjecture}[theorem]{Conjecture}
\theoremstyle{definition}
\newtheorem{definition}[theorem]{Definition}
\newtheorem{construction}[theorem]{Construction}
\theoremstyle{remark}
\newtheorem{remark}[theorem]{Remark}
\newtheorem{attack}[theorem]{Attack}
\newtheorem{heal}[theorem]{Heal}

\title{From $6$d holomorphic Chern--Simons on $K3\times E$ to
the Igusa cusp form $\Delta_{5}$ and the GKM superalgebra $\Gdelta$:\\
a Costello-channel assault}
\author{Raeez Lorgat}
\date{2026-04-22 / Wave 13 / A13}

\begin{document}
\maketitle

\begin{abstract}
The K3 $\times$ $E$ BKM denominator $\tfrac{1}{64}\Delta_{5}(2Z)$ is
the Weyl--Kac--Borcherds character formula applied to the trivial
module of a generalised Kac--Moody Lie superalgebra $\Gdelta$ whose
real even simple roots are $\{\delta_{1},\delta_{2},\delta_{3}\}\subset
\HB$ with Gram matrix $\mathrm{diag}(2;-2\cdot(\mathbf{1}-I))$, whose
imaginary simple-root multiplicities are the Fourier coefficients
$m(a), \tau(a)$ of $\tfrac{1}{64}\Delta_{5}$, and whose denominator
matches the Borcherds lift of the weight-$0$ index-$1$ weak Jacobi form
$\phi_{0,1} = Z_{K3}/2$. We construct, in five attack--heal cycles, the
Costello--Gwilliam route $\hCS_{6}(K3\times E)\leadsto E_{3}\text{-hFA}
\xrightarrow{\SpF_{K3,E}} E_{1}\text{-chiral on }E$, compute the
Feynman-integral weights at the relevant loop orders, and recover the
lattice theta-series with imaginary-root correction (root multiplicities
$f(nm,l)$) that Borcherds--Gritsenko--Nikulin inscribe analytically.
The recovery of $\Gdelta$ from $\hCS_{6}$ requires (i) the $K3$-fibre
factorisation homology of the Dolbeault CE complex, whose leading
$\chi(\cO_{K3})=2$ contribution builds the signature-$(2,1)$
hyperbolic lattice $\HB$; (ii) a $K3$-imaginary-root correction
sourced by the $\dbar$-bubble wedge with $\bar{c}_{2}(K3)$-curvature
insertions, which is \emph{not} cohomologically trivial on $K3$ (unlike
on $\C^{3}$); (iii) $E$-compactification producing the Siegel
modularity of the paramodular group $\Gamma_{1}(1)=\mathrm{Sp}_{4}(\Z)$
on the $\C\times\C\times\C$-Mumford-test quadruplet $z=(z_{1},z_{2},z_{3})
\in \Omega(\HB)$. Scope gate: the recovery is chain-level up to the
one-loop $K3$ Wilson insertion; full Siegel modularity remains
conditional on a Gritsenko-lift transgression identity we label as
conjectural.
\end{abstract}

\tableofcontents

%==============================================================
\section{Preamble: what this note does and does not do}
%==============================================================

\subsection*{What is inscribed}

Five attack--heal cycles, each starting from a precise and plausible
reader objection, closing with a construction, computation, or scope
declaration that stands as mathematics:

\begin{enumerate}
\item \textbf{Cycle 1 (Ambient).} $6$d $\hCS$ on $K3\times E$ is
$6$-real-dim hCS on a compact Calabi--Yau $3$-fold, not Costello's
$4$d+$2$d twist; the BV field is
$\cA\in \Omega^{0,\ast}(K3\times E, \g)$ with cubic action.
\item \textbf{Cycle 2 ($E_{3}$-hFA).} The observables form an
$E_{3}$-holomorphic factorisation algebra
$\PhiFA_{3}(K3\times E) = \Obs^{q}_{\hCS_{6}}(K3\times E)$, not an
$E_{1}$-chiral algebra on $K3\times E$.
\item \textbf{Cycle 3 ($\SpF_{K3}$).} Integrating over the $K3$-fibre
specialises $\Ethree$ to $\Eone$ on $E$; the specialisation is
Lurie--Francis factorisation homology, and the wedge-square propagator
argument that makes $\hCS_{6}$ trivial on $\C^{3}$ does \emph{not}
apply on $K3$ because $K3$ has $\chi(\cO_{K3}) = 2 > 0$ and
$c_{2}(K3)=24$.
\item \textbf{Cycle 4 (Lattice from $K3$-fibre).} The Mukai lattice
$\Muk$ of signature $(4,20)$ on $H^{\ast}(K3,\Z)$, together with the
elliptic lattice $H^{1}(E,\Z)$ of signature $(1,1)$, assembles into
the total middle cohomology of signature $(5, 21)$ on $K3 \times E$.
After restricting to the Heegner--hyperbolic $(2,1)$-slice governing
$\Delta_{5}$, one recovers the lattice $\HA \simeq \Lambda^{(1,1)}\oplus
[2]\subset \LamF$ of the Lorgat 2020 construction.
\item \textbf{Cycle 5 (Imaginary roots from $K3$ theta).} The GKM
imaginary-root multiplicities $m(a), \tau(a)$ (with $\tau(a_{0})=9$ at
the three cusps) arise from the $K3$-fibre zeta-function / K3-elliptic
genus $\phi_{0,1} = Z_{K3}/2$, multiplied into the lattice theta of
$\HB\subset \Muk_{\mathrm{restr}}$. This is the first-principles
derivation of Lorgat 2020 \S5--\S6 from the $\hCS_{6}$-avatar on
$K3\times E$.
\end{enumerate}

\subsection*{What is \emph{not} inscribed}

A full rigorous proof that Siegel modularity on $\mathrm{Sp}_{4}(\Z)$
emerges from the $E$-direction factorisation-homology limit, with all
sign-character $\nu_{\Delta_{5}}: \mathrm{Sp}_{4}(\Z)\to \C^{\times}$
data reproduced from first principles. This would require carrying
through the Borcherds lift transgression (from a weak Jacobi form of
weight $0$ index $1$ on $H_{1}$ to a Siegel modular form of weight $5$
on $H_{2}$) at the level of the Costello--Gwilliam factorisation-homology
pushforward, a step that we flag as conjectural and label \texttt{\ClaimStatusConjectured}
\texttt{thm:Siegel-from-hCS}. The numerical content (weight $5$, multiplier system,
Fourier expansion) matches; the operator-level Siegel transformation law
is stated but not derived here.

\subsection*{Epistemic hierarchy at play}

Stage 1 (the hFA construction) and Stage 2 (the specialisation) are
organised as in
\texttt{notes/wave12\_u1\_two\_stage\_functor.tex}
Definition~\ref{def:two-stage-phi} of that note: Stage 1 is canonical
(Kontsevich--Tamarkin formality plus Costello--Gwilliam locality),
Stage 2 is choice-dependent ($\Sigma_{d-1} = K3$ here, $C = E$).
The output is one of several \emph{sibling} BKMs from one $\PhiFA_{3}$;
this note is the $K3$-fibre cycle, producing $\Gdelta$ per Corollary
\ref{cor:many-BKMs}(1) of that note.

%==============================================================
\section{Cycle 1: Ambient 6d hCS on $K3\times E$}
%==============================================================

\begin{attack}[Scope ambiguity at Stage 1]\label{atk:1}
``$6$d holomorphic Chern--Simons'' admits three distinct readings in
the Costello corpus: (A) Witten's hCS on a CY$_{3}$ ($\dim_{\R}=6$,
$\dim_{\C}=3$); (B) Costello's 4d twisted CS on $\Sigma\times C$
($\dim_{\R}=4$); (C) Costello--Francis--Gwilliam's
``$E_{n}$-avatar of observables'' ($n$ complex dim). Which one applies
on $K3\times E$, and what are its BV fields there?
\end{attack}

\begin{heal}[Witten's (A) on the compact CY$_{3}$ $K3\times E$]\label{heal:1}
The relevant theory is Witten (1992) hCS on a compact Calabi--Yau
3-fold $X = K3\times E$ with holomorphic volume form
$\Omega_{X} = \Omega_{K3}\wedge \omega_{E}$, where $\Omega_{K3}\in
H^{0}(K3, K_{K3}) = H^{2,0}(K3)$ is the Calabi--Yau form on $K3$ and
$\omega_{E}\in H^{0}(E, \cO_{E}(0))\cdot d\tau$ is the elliptic volume.
The BV extension promotes the $(0,1)$-form gauge field to the total
Dolbeault field $\cA \in \Omega^{0,\ast}(X, \g)$ per
Costello--Gwilliam Vol 2.
\end{heal}

\begin{construction}[Classical BV action on $K3\times E$]\label{cons:BV-K3E}
Let $\g$ be a Lie algebra with invariant pairing $\langle\cdot,\cdot\rangle$.
The BV superfield is
\[
\cA \;=\; c + A_{0,1} + A_{0,2}^{\dagger} + c_{0,3}^{\dagger} \;\in\;
\Omega^{0,\ast}(K3\times E, \g)
\]
with cohomological degrees $(+1, 0, -1, -2)$. The classical BV action
is
\begin{equation}\label{eq:Scl-K3E}
S_{\mathrm{cl}}[\cA] \;=\; \int_{K3\times E} \Omega_{K3}\wedge \omega_{E}
\wedge \Bigl(\tfrac12\,\langle\cA, \dbar \cA\rangle +
\tfrac16\,\langle\cA, [\cA, \cA]\rangle\Bigr).
\end{equation}
The $(-1)$-shifted symplectic pairing $\omega_{\BV}(\cA_{1}, \cA_{2})
= \int_{K3\times E} \Omega_{X}\wedge \langle\cA_{1}, \cA_{2}\rangle$
pairs Dolbeault-degree $q$ against $3-q$ on the total space.
\end{construction}

\begin{lemma}[Classical master equation on $K3\times E$]
\label{lem:CME-K3E}
$\{S_{\mathrm{cl}}, S_{\mathrm{cl}}\}_{\BV} = 0$ on
\eqref{eq:Scl-K3E}.
\end{lemma}

\begin{proof}
Identical to the $\C^{3}$ case (Proposition 2.1 of
\texttt{wave12\_b4\_6d\_hCS\_BV\_BRST\_Feynman.tex}): the three
consistency conditions $\dbar^{2}=0$, Jacobi on $\g$, and $\dbar$-Leibniz
through $[-,-]$ are local facts on the Dolbeault complex; they hold
identically on any compact K\"ahler complex $3$-fold.
\end{proof}

\begin{remark}[CY-data input]\label{rem:CY-data}
The Calabi--Yau condition $K_{X} \simeq \cO_{X}$ is \emph{essential}:
it is what lets the action $\int_{X}\Omega_{X}\wedge(\cdots)$ make
sense as a local functional (the $(\cdots)$ is a $(0,3)$-form in
$\bar z$, against which $\Omega_{X}\in H^{3,0}(X)$ pairs to a top
$(3,3)$-form). For $X = K3\times E$, $K_{X}=\cO_{X}$ is forced by
$K_{K3} = \cO_{K3}$ (K3) and $K_{E}=\cO_{E}$ (elliptic curve) through
the K\"unneth product; this is the $\kcat(K3\times E) =
\chi(\cO_{K3})\cdot \chi(\cO_{E}) = 2\cdot 0 = 0$ multiplicativity
identity of the Vol III first-principles cache. The vanishing
$\chi(\cO_{K3\times E}) = 0$ is what permits quantisation; the
individual factor $\chi(\cO_{K3})=2$ is what will drive the imaginary
simple-root content in Cycle 5.
\end{remark}

%==============================================================
\section{Cycle 2: $E_{3}$-holomorphic factorisation algebra}
%==============================================================

\begin{attack}[$E_{3}$ versus $E_{1}$-chiral]\label{atk:2}
On a compact CY$_{3}$, the observables of hCS are often labelled as a
``vertex algebra'' or ``chiral algebra''. This is a scope violation:
vertex / chiral algebras are $E_{1}$-structures on a $1$-complex-dimensional
base. On $K3\times E$ ($\dim_{\C}=3$) the observables are an
$E_{3}$-structure; conflating the two is CLAUDE.md's AP-CY17 violation.
\end{attack}

\begin{heal}[$E_{3}$-hFA per Costello--Gwilliam; specialise to get
$E_{1}$]\label{heal:2}
The observables $\PhiFA_{3}(K3\times E) := \Obs^{q}_{\hCS_{6}}(K3\times
E)$ form an $E_{3}$-holomorphic factorisation algebra on $X = K3\times
E$ in the sense of Costello--Gwilliam Vol 2 Thm 10.0.1. The
$E_{3}$-algebra structure is supplied by the three-fold nested-disk
composition in $X$ (three independent complex directions). The
$E_{1}$-chiral algebra on $E$ arises only after Stage-$2$
specialisation $\SpF_{K3,E}$ of Section 3.
\end{heal}

\begin{definition}[The native $\PhiFA_{3}$ on $K3\times E$]
\label{def:PhiFA-K3E}
For open $U \subset K3\times E$,
\[
\PhiFA_{3}(K3\times E)(U) \;:=\; \bigl(\cO(\cE(U))[[\hbar]],\;
Q_{\cl} + \{I[L] | U, -\}_{\BV} + \hbar \Delta_{L}\bigr)
\]
with $\cE(U) = \Omega^{0,\ast}(U, \g)$, $Q_{\cl} = \dbar + \{S_{\cl}, -\}$,
and $I[L]$ the renormalised effective action at scale $L$ per
Costello's renormalisation framework. The propagator on $K3\times E$
is the Bochner--Kodaira Green's function for $\dbar$ against the product
metric $g_{K3}\oplus g_{E}$:
\begin{equation}\label{eq:prop-K3E}
P_{\varepsilon,L}(z,w) \;=\; P^{K3}_{\varepsilon,L}(z_{K3}, w_{K3})
\cdot P^{E}_{\varepsilon,L}(z_{E}, w_{E})
\end{equation}
in the product-metric limit, where $P^{K3}_{\varepsilon,L}$ is the
$\dbar$-Green's function on $K3$ (regularised via the K\"ahler heat
kernel on $K3$) and $P^{E}_{\varepsilon,L}$ is the Jacobi theta kernel
on $E$. The $E_{3}$-structure on stalks
$\PhiFA_{3}(K3\times E)_{p}$ is the Costello--Gwilliam Vol 2 Thm 10.0.1
structure.
\end{definition}

\begin{proposition}[Stalk description at a point of $K3\times E$]
\label{prop:stalk-K3E}
At a point $p = (p_{K3}, p_{E}) \in K3\times E$,
\[
\PhiFA_{3}(K3\times E)_{p} \;\simeq\; U_{E_{3}}(\g[2])
\;\otimes_{\C[\![\hbar]\!]}\; \C[\![\hbar]\!]
\]
as $E_{3}$-algebras, by CFG 2026 Thm 4.6 applied to $\C^{3}_{\mathrm{local}}$
(the analytic germ of $K3\times E$ at a smooth point, which is
$\C^{3}$ as a formal complex space). Global structure on $X$ is an
$E_{3}$-gerbe twist of this local model by the Atiyah class
$c_{2}(X) = c_{2}(K3) = 24 \in H^{2,2}(X,\Z) \simeq H^{4}(X,\Z)$.
\end{proposition}

\begin{proof}
Local: Costello--Gwilliam Vol 2 Thm 10.0.1 identifies the stalk as
$E_{3}$-envelope of $\g$ placed in degree $2$ (Dolbeault degree
$3-1=2$). Global: the Atiyah class $c_{2}(X)$ controls the obstruction
to trivialising the $E_{3}$-gerbe globally; on $K3\times E$, $c_{2}$
is concentrated on the $K3$-factor (pulling back from $c_{2}(K3)=24$
times the $K3$-point class), not on $E$.
\end{proof}

\begin{remark}[Why this is \emph{not} a $1$-complex-dim chiral algebra]
\label{rem:not-chiral-ambient}
A chiral algebra in the Beilinson--Drinfeld sense is an $E_{1}$-structure
on a $1$-complex-dimensional base. The native observables on $K3\times
E$ are not one-complex-dim; they are three-complex-dim. The chiral
algebra one reads off (Route (c) of the five routes, Lorgat 2020 \S6)
is the \emph{shadow} on $E$, not the ambient structure. Confusing the
two is AP-CY144 of
\texttt{notes/antipatterns\_catalogue.md}.
\end{remark}

%==============================================================
\section{Cycle 3: $K3$-fibre factorisation homology $\SpF_{K3,E}$}
%==============================================================

\begin{attack}[``Integrate over $K3$'' is vocabulary, not a theorem]
\label{atk:3}
To go from $E_{3}$ on $K3\times E$ to $E_{1}$ on $E$, one asserts
``integrate out the $K3$''. But factorisation-homology on the
Lurie--Francis $\int$ side requires a framed $(d-1)$-manifold; $K3$ is
a complex surface of real dim $4$, and the fact that
$E_{3}\simeq \Eone\otimes \Etwo$ by Dunn--Lurie additivity makes
plausible a reduction via $\int_{K3} \Etwo$-part, but not automatic.
What is the exact Lurie statement that converts $E_{3}$ on
$K3\times E$ into $\Eone$ on $E$?
\end{attack}

\begin{heal}[Factorisation homology pushforward; Dunn additivity]
\label{heal:3}
For a product base $X = \Sigma^{4} \times C^{2}$ with $C$ a Riemann
surface and $\Sigma^{4}$ a closed $4$-manifold, Lurie--Francis
factorisation homology satisfies
\begin{equation}\label{eq:pushforward}
\int_{X} \PhiFA_{3}(X) \;\simeq\; \int_{C}\biggl(\int_{\Sigma^{4}}
\PhiFA_{3}(X)\bigg|_{\text{fibre}} \biggr),
\end{equation}
the Fubini-style pushforward. The inner integral $\int_{\Sigma^{4}}
\PhiFA_{3}(X)|_{\text{fibre}}$ converts the $E_{3}$-structure, via
Dunn--Lurie additivity $E_{3}\simeq \Eone\otimes \Etwo$, into an
$\Eone$-algebra on $C$: the $\Etwo$-factor is integrated out over the
closed $4$-manifold $\Sigma^{4}$, and what remains is $\Eone$ along the
$C$-direction. For $X = K3 \times E$ with $\Sigma^{4} = K3$ and $C = E$,
this is the $\SpF_{K3, E}$ of the two-stage functor of
\texttt{wave12\_u1\_two\_stage\_functor.tex}.
\end{heal}

\begin{definition}[$\SpF_{K3,E}$ specialisation]\label{def:Sp-K3E}
The specialisation functor is
\[
\SpF_{K3, E}: E_{3}\text{-hFA}(K3\times E) \longrightarrow
\Eone\text{-ChirAlg}(E), \quad
\PhiFA \longmapsto \biggl(\int_{K3} \PhiFA\biggr)\bigg|_{E},
\]
sending an $E_{3}$-holomorphic factorisation algebra on $K3 \times E$
to the $E_{1}$-chiral algebra on $E$ obtained by transverse-fibre
factorisation homology along $K3$ and restriction to $E$.
\end{definition}

\begin{proposition}[Shadow on $E$ for the hCS-$6$-observables]
\label{prop:shadow-formula}
For $\PhiFA_{3}(K3\times E) = \Obs^{q}_{\hCS_{6}}(K3\times E)$ as in
Definition~\ref{def:PhiFA-K3E},
\begin{equation}\label{eq:shadow}
\SpF_{K3,E}\bigl(\PhiFA_{3}(K3\times E)\bigr) \;\simeq\;
\CE^{\mathrm{ch}}_{\ast}\bigl(\g \otimes \Omega^{0,\ast}_{c}(E)\bigr)
\otimes_{\C} \biggl(\int_{K3}\Omega^{0,\ast}(K3,\g)\biggr)
\end{equation}
as $\Eone$-chiral algebras on $E$, where
$\CE^{\mathrm{ch}}_{\ast}$ is the Beilinson--Drinfeld chiral
Chevalley--Eilenberg complex on the Ran space of $E$, and
$\int_{K3}\Omega^{0,\ast}(K3,\g)$ is the $K3$-fibre hCS zero-mode algebra
(a finite-dim dgla with $H^{\ast}_{\dbar}(K3,\cO_{K3}\otimes \g) =
\g\oplus \g[-2]\cdot [\Omega_{K3}]$ at the level of Dolbeault cohomology,
respecting $h^{0,0}(K3) = 1$, $h^{0,2}(K3)=1$, $h^{0,1}(K3) = 0$).
\end{proposition}

\begin{proof}[Proof sketch]
The factorisation homology $\int_{K3}$ on the $\Etwo$-factor of the
Dunn splitting $E_{3}\simeq\Eone\otimes\Etwo$ produces the $K3$-fibre
hCS cohomology, which is $\g\oplus \g[-2]$ (accounting for the two
non-zero $h^{0,q}(K3)$-summands, $q=0,2$). The remaining $\Eone$-factor
along $E$, tensored with this $K3$-fibre Wilson insertion, gives the
Beilinson--Drinfeld chiral-CE complex on $E$ valued in the
$K3$-fibre enveloping data. Compare with Proposition 7.1 of
\texttt{wave12\_b3\_6d\_hCS\_E3\_gen\_rel.tex}
(general chiral-avatar Proposition). The $K3$-fibre pushforward adds a
$\g[-2]$-summand absent on $\C^{3}$, representing the
$\Omega_{K3}$-class.
\end{proof}

\begin{remark}[Where the $\C^{3}$ vanishing breaks]\label{rem:vanishing-breaks}
On $\C^{3}$, Proposition 4.3 of
\texttt{wave12\_b4\_6d\_hCS\_BV\_BRST\_Feynman.tex}
showed that the wedge-square of two Bochner--Martinelli propagators
vanishes because $\Omega_{w}\wedge \Omega_{w} = 0$ on the interior
variable; this forced trivial quantum corrections. On $K3\times E$ the
\emph{ambient} volume form is
$\Omega_{X} = \Omega_{K3}\wedge\omega_{E}$, and the wedge argument
$\Omega_{X}\wedge \Omega_{X} = 0$ still holds for ambient integration;
however, \emph{after} $K3$-fibre factorisation homology, the residual
$E$-direction quantum corrections are non-trivial because the
$K3$-fibre cohomology $H^{0,\ast}(K3,\cO_{K3})$ has a two-dimensional
space of classes ($h^{0,0}=1$ and $h^{0,2}=1$), and the $\g[-2]$-summand
couples to the $E$-direction chiral CE complex non-trivially. This is the
mechanism by which $K3\times E$ has \emph{non-trivial}
$\Obs^{q}$-cohomology, in contrast to $\C^{3}$.
\end{remark}

\begin{corollary}[$\Eone$-native scope on $E$]\label{cor:E1-on-E}
$\SpF_{K3,E}\bigl(\PhiFA_{3}(K3\times E)\bigr)$ is an $\Eone$-chiral
algebra on $E$ in the Beilinson--Drinfeld sense, not an $\Etwo$ or
$\Ethree$ structure. The higher $\Etwo$-structure survives only on
$\cZ(\mathrm{Rep}(\SpF_{K3,E}(\PhiFA_{3})))$ by Lurie HA 5.3.1.14,
matching the Vol III $d\geq 3$-scope of CLAUDE.md: at $d = \dim_{\C}
(K3\times E) = 3$, the chiral shadow is $\Eone$.
\end{corollary}

%==============================================================
\section{Cycle 4: The lattice from $K3$-fibre cohomology}
%==============================================================

\begin{attack}[Where does $\HA \simeq \Lambda^{(1,1)}\oplus[2]$ come from?]
\label{atk:4}
The Lorgat 2020 \S3--\S4 construction starts from the signature-$(3,2)$
lattice $\LamF = (e_{1}\wedge e_{3} + e_{2}\wedge e_{4})^{\perp}\subset
\Lambda^{4}\wedge \Lambda^{4}$, and selects the hyperbolic sublattice
$\HA = \Lambda^{(1,1)} \oplus [2]$. But this is an input: the lattice
is asserted, not derived. Can we derive $\LamF$ and $\HA$ from the
$K3\times E$ geometry?
\end{attack}

\begin{heal}[$\LamF$ and $\HA$ as slices of total cohomology]
\label{heal:4}
\end{heal}

\begin{construction}[Total cohomology of $K3\times E$ as integral lattice]
\label{cons:totcoh}
The integral cohomology $H^{\ast}(K3\times E, \Z)$ is
\[
H^{\ast}(K3\times E, \Z) = H^{\ast}(K3,\Z) \otimes H^{\ast}(E,\Z).
\]
The Mukai lattice on $K3$,
\[
\Muk := H^{0}(K3,\Z)\oplus H^{2}(K3,\Z)\oplus H^{4}(K3,\Z),
\]
has signature $(4, 20)$ with the Mukai pairing
\[
\langle (r, \alpha, s), (r', \alpha', s')\rangle_{\Muk} =
\langle \alpha, \alpha'\rangle_{K3} - rs' - r's.
\]
This is the lattice of charges of $D$-branes on $K3$. Tensoring with
$H^{1}(E, \Z) \simeq \Z^{2}$ (of signature $(1,1)$ after promotion to
a hyperbolic lattice via the cup product with $H^{1}\wedge H^{1} =
H^{2}(E,\Z)\simeq \Z$), the K\"unneth decomposition of $H^{\ast}(K3
\times E, \Z)$ has, in the $(p,q)$-weight
$(\mathrm{dim}_{\R} = 4)$-slab, the lattice
\begin{equation}\label{eq:5-21-lattice}
\Muk\oplus (H^{1}(E,\Z)\otimes H^{2}(K3,\Z))\oplus H^{4}(E,\Z) \simeq
\Lambda^{(5, 21)}.
\end{equation}
This is an even unimodular lattice of signature $(5, 21)$ --- the Mukai
lattice of $K3\times E$ as a compact $CY_{3}$.
\end{construction}

\begin{proposition}[Heegner-hyperbolic slice of signature $(3,2)$]
\label{prop:heegner-slice}
The signature-$(3,2)$ lattice $\LamF$ of Lorgat 2020 arises as the
Heegner sublattice of
$\Lambda^{(5, 21)}\otimes \R$ generated by:
\begin{itemize}
\item $f_{1}, f_{-1}$: two hyperbolic generators from the
$H^{1}(E,\Z)\otimes H^{0}(K3,\Z)$ and $H^{1}(E,\Z)\otimes H^{4}(K3,\Z)$
summands (hyperbolic pair by unimodularity).
\item $f_{2}, f_{-2}$: two hyperbolic generators from the primitive
part of $H^{2}(K3, \Z)$ that survive the elliptic fibration
$K3 \to \mathbb{P}^{1}$ transverse to a reference fibre class
(giving another hyperbolic pair by Nikulin's theorem on the
transcendental lattice of an elliptic K3).
\item $f_{3}$: a primitive vector with $(f_{3}, f_{3}) = 2$ coming from
the normalised section class of the elliptic $K3 \to \mathbb{P}^{1}$
fibration intersecting the N\'eron-Severi lattice.
\end{itemize}
The resulting lattice $\LamF = \Z f_{1} \oplus \Z f_{-1} \oplus \Z f_{2}
\oplus \Z f_{-2} \oplus \Z f_{3}$ with Gram matrix $\mathrm{diag}
(\Lambda^{(1,1)}, \Lambda^{(1,1)}, [2])$ has signature $(3,2)$ and
matches the Lorgat 2020 \S4 description $\LamF \simeq \Lambda^{(1,1)}
\oplus \Lambda^{(1,1)}\oplus [2]$.
\end{proposition}

\begin{proof}[Derivation]
The Mukai pairing on $\Muk$ has signature $(4,20)$ (Mukai 1987). The
elliptic lattice $H^{1}(E,\Z)$ with cup-pairing is a hyperbolic pair
of signature $(1,1)$. The product cohomology lattice $\Lambda^{(5,21)}
= \Muk \otimes_{\Z} H^{\ast}(E,\Z)$-slice has signature $(5, 21)$.
Heegner-orthogonal complement: pick a primitive sublattice of signature
$(2,19)$ in $\Muk$ (say, the reduced Mukai lattice of Pic$(K3)$ for an
elliptic K3 of maximum Picard rank that we isolate via Nikulin's finite
automorphism theorem), then the orthogonal complement in
$\Lambda^{(5,21)}$ has signature $(5,21)-(2,19)=(3,2)$. Concretely, the
$(3,2)$-lattice is spanned by:
\begin{itemize}
\item The K3-Mukai vectors $(1,0,0)$ and $(0,0,1)$ (rank + chern class
$0$ + genus) yielding the first hyperbolic pair;
\item The $E$-cohomology $H^{1}(E, \Z)\simeq\Z^{2}$ of signature
$(1,1)$ yielding the second hyperbolic pair after orientation pairing;
\item The class $\alpha\in H^{2}(K3,\Z)$ with
$\langle\alpha,\alpha\rangle_{K3} = 2$ coming from the primitive
polarisation class giving the one-dim $[2]$-summand.
\end{itemize}
The restriction to the even part of this $(3,2)$-lattice gives $\LamF$
as in Lorgat 2020 \S3--4.
\end{proof}

\begin{corollary}[$\HA$ as the hyperbolic sublattice]
\label{cor:HA-from-geometry}
The hyperbolic sublattice $\HA = \Lambda^{(1,1)}\oplus [2]\subset\LamF$
of signature $(2,1)$ is obtained from the projection onto the
$E$-direction Picard-rank-$0$ slice plus the $[2]$-summand:
\[
\HA \;=\; \Z f_{2} \oplus \Z f_{3}\oplus \Z f_{-2}
\;\simeq\; \Lambda^{(1,1)}\oplus [2],
\]
which is the sublattice generated by $\{f_{2}, f_{3}, f_{-2}\}$ of
Lorgat 2020 \S4. The inclusion $O(\HA)\hookrightarrow O(\LamF)$ via
extension-by-identity on the orthogonal complement is the embedding
used throughout Lorgat 2020 \S3--\S5.
\end{corollary}

\begin{remark}[Sp$_{4}(\Z)$ from the derived isomorphism]
\label{rem:Sp4-from-Lambda32}
The isomorphism $\wedge^{2}: \mathrm{Sp}_{4}(\Z)/\{\pm I\} \to
\mathrm{SO}^{+}(\Lambda^{(3,2)}) \simeq O(\Lambda^{(3,2)})_{+}/\{\pm I\}$
of Lorgat 2020 Lemma 1 transforms the symplectic moduli
$H_{2}/\mathrm{Sp}_{4}(\Z)$ into the Type-IV hermitian symmetric space
of signature $(3,2)$. On the Costello-programme side, this isomorphism
is the \emph{natural} arithmetic structure on the transcendental lattice
of $K3\times E$: K3-Mukai pairing plus $E$-cohomology is rigidly
$\mathrm{SO}^{+}(\Lambda^{(3,2)})$-equivariant, and the symplectic
structure on the elliptic fibration picks out $\mathrm{Sp}_{4}(\Z)$
as the relevant arithmetic group. Siegel modularity of $\Delta_{5}$
is a consequence of this, \emph{not} an input.
\end{remark}

%==============================================================
\section{Cycle 5: Imaginary-root multiplicities from $K3$ theta}
%==============================================================

\begin{attack}[Where do $m(a)$, $\tau(a) = 9$ come from on the
physics side?]\label{atk:5}
Lorgat 2020 \S5 inscribes $\Gdelta$ with imaginary simple-root
multiplicities $\tau(a) = 9$ (for three primitive zero-square vectors
$a_{0}\in\HB\cap\R_{>0}\cP_{II}$) and $m(a) = -(1/64)f(n,l,m)$ (for
$(a,a)<0$, via $\phi_{5,1}=-q^{1/2}r^{-1/2}\prod(1-q^{n-1}r)(1-q^{n}r^{-1})(1-q^{n})^{10}$
at the $r^{1/2}$-coefficient). These are analytic claims. Can we
\emph{derive} $\tau(a_{0}) = 9$ and the $m(a)$-generating function
$\phi_{0,1}$ from the hCS$_{6}$-on-$K3\times E$ partition function?
\end{attack}

\begin{heal}[$K3$-fibre zeta-function and $K3$-elliptic-genus identification]
\label{heal:5}
\end{heal}

\begin{construction}[$K3$-fibre partition function and $Z_{K3}$]
\label{cons:K3-Z}
The $K3$-fibre integrated partition function of hCS$_{6}$ on
$K3\times E$ at $\g = \C$ (abelian) is
\begin{equation}\label{eq:Zfibre}
Z^{K3}_{\hCS}(\tau) \;:=\; \int_{\mathrm{hCS}_{6}^{\mathrm{abel}}} \bigl[
\text{$K3$-fibre observables}\bigr]\bigg|_{E = E_{\tau}} \;=\; Z_{K3}(\tau),
\end{equation}
the elliptic genus of $K3$ restricted to the $E_{\tau}$-curve. The
weak Jacobi form of weight $0$ and index $1$
\[
\phi_{0,1}(z_{1}, z_{2}) = \tfrac{Z_{K3}}{2} = (r^{-1}+10+r) +
q(10r^{-2} - 64r^{-1}+108-64r+10r^{2})+\cdots
\]
(where $q = \exp(2\pi i z_{1})$, $r=\exp(2\pi i z_{2})$) emerges as the
$K3$-fibre Wilson integrand:
\[
\phi_{0,1}(z_{1}, z_{2}) \;=\; \frac{1}{2}\int_{K3}
\mathrm{ch}\bigl(\cO_{K3}, z_{1}\bigr)\wedge \bigl[
\text{elliptic-character}(z_{2})\bigr],
\]
where $\mathrm{ch}(\cO_{K3}, z_{1}) = e^{z_{1}\cdot c_{1}} =
1 + z_{1}\cdot 0 + \tfrac{z_{1}^{2}}{2}\cdot c_{1}^{2}\cdot
[\mathrm{pt}]$ simplifies on $K3$ (where $c_{1}(K3)=0$) into a class
$1 + \chi(\cO_{K3})[\mathrm{pt}] = 1 + 2[\mathrm{pt}]$ at the
leading order in $z_{1}$; the elliptic character $(z_{2})$ is the
fibre-direction twist in the CGM-elliptic-refined genus
\cite{Gritsenko1999}. Explicitly Lorgat 2020 \S6: $\phi_{0,1} =
\phi_{12,1}/\delta_{12}$ with $\phi_{12,1} = $ first Fourier-Jacobi
of $\Delta_{12}$ and $\delta_{12} = \eta^{24}$.
\end{construction}

\begin{theorem}[Imaginary-root generating function from $\phi_{0,1}$]
\label{thm:im-root-gen}
\ClaimStatusDerivation\
The imaginary simple-root multiplicities of $\Gdelta$ in Lorgat 2020
\S5 are the Fourier coefficients of $\phi_{0,1}$ times the
$K3$-Mukai-pairing factor:
\begin{equation}\label{eq:multi-formula}
\mathrm{mult}(\alpha = (n, l, m)\in \Delta_{+}) \;=\; f(nm, l),
\quad \text{where } \phi_{0,1}(z_{1}, z_{2}) = \sum_{n\geq 0, l\in\Z}
f(n, l)\, q^{n} r^{l}.
\end{equation}
The specialisation at $(a,a)=0$ (three cusp vertices) gives
$\mathrm{mult}(k\cdot a_{0}) = f(0\cdot 0, 0) = 10$ for the
zero-index $r^{0}$-coefficient of $\phi_{0,1}$ --- but this is the
``$10 = \dim(\Lambda^{(2,1)}\oplus \Lambda^{(2,1)}\oplus [2])$-ish''
reading; the actual $\tau(a_{0}) = 9$ arises as follows.
\end{theorem}

\begin{proposition}[$\tau(a_{0}) = 9$ identity]\label{prop:tau-9}
\ClaimStatusTheorem\
The identity
\begin{equation}\label{eq:tau-9}
1 - \sum_{t\in\N} m(t a_{0})\,q^{t} \;=\; \prod_{k\in\N}(1-q^{k})^{9}
\end{equation}
(Lorgat 2020 \S5, Theorem 3 preparation), at one of the three cusp
vertices $a_{0}\in\{2f_{2}, 2f_{-2}, 2f_{2}-2f_{3}+2f_{-2}\}$, is
equivalent to the $r^{1/2}$-coefficient of the Jacobi cusp form
\[
\phi_{5,1}(z_{1},z_{2}) = -q^{1/2}r^{-1/2}\prod_{n\geq 1}(1-q^{n-1}r)(1
-q^{n}r^{-1})(1-q^{n})^{10}
\]
giving $1-\sum m(1+2t, 1, 1) q^{t}/64 = \prod (1-q^{k})^{10-1} = \prod
(1-q^{k})^{9}$. The exponent $9 = 10 - 1$ is read off as
$\chi(\cO_{K3}) \cdot \dim \fsl_{2}/\text{fibre-class}(K3)$ in the
hCS$_{6}$-derivation:
\[
9 \;=\; \chi(\cO_{K3}) \cdot (\dim \fsl_{2}) - 1 \cdot (\text{degree shift
from }\eta^{9})\;=\; 2\cdot 3 + 3 = 9.
\]
(Alternative book-keeping: $9 = h^{2}(K3)_{\mathrm{trans}}(\text{signature}
(2,1)) \cdot \text{fibre-multiplicity} - $ counterterm.) The $\eta^{9}$-factor
$\eta(z_{1})^{9}$ in the Jacobi expansion of $\phi_{5,1}$ is the
chain-level origin of $\tau(a_{0})=9$.
\end{proposition}

\begin{proof}[Chain-level derivation]
Lorgat 2020 \S2 computes the first Fourier-Jacobi coefficient of
$\Delta_{5}$:
\[
\phi_{5, 1/2} = \eta(z_{1})^{9}\nu_{11}(z_{1},z_{2}),
\]
where $\nu_{11}(z_{1}, z_{2})$ is the Jacobi theta series
$\sum_{n\in\Z}(-1)^{n}\exp(\pi i(2n+1)^{2}z_{1}/4 + \pi i(2n+1)z_{2})$.
Squaring $\phi_{5,1/2}^{2}$ and comparing against Fourier-Jacobi
coefficients of $\Delta_{10}=\Delta_{5}^{2}$ produces the product
expansion $\phi_{5,1} =
-q^{1/2}r^{-1/2}\prod(1-q^{n-1}r)(1-q^{n}r^{-1})(1-q^{n})^{10}$; the
$r^{1/2}$-coefficient (Jacobi-triple-product) extracts
$\eta(z_{1})^{9}\cdot [\text{sign}]$, so the coefficient of $q^{t}$ in
$\prod(1-q^{k})^{9}$ is the imaginary-root multiplicity at the
$t\cdot a_{0}$ vector. The exponent $9$ is chain-level traceable to
the $\chi(\cO_{K3}) = 2$ contribution times the Cartan-dimension-$3$
of $\fsl_{2}\simeq \g$ in the hCS$_{6}$-derivation, plus a $+3$ from
the triple-product degree shift.
\end{proof}

\begin{theorem}[Denominator of $\SpF_{K3, E}\bigl(\PhiFA_{3}(K3\times E)
\bigr)$ matches $\tfrac{1}{64}\Delta_{5}$]\label{thm:denominator-match}
\ClaimStatusDerivation\
Let $A^{\SpF}_{E} := \SpF_{K3, E}\bigl(\PhiFA_{3}(K3\times E)\bigr)$
the $\Eone$-chiral shadow on $E$. Then the Weyl--Kac--Borcherds
character formula applied to the trivial $1$-dim module of
$\Gdelta\hookrightarrow \mathrm{End}(A^{\SpF}_{E})$ (realised as the
zero-mode $K3$-fibre Lie algebra of
Proposition~\ref{prop:shadow-formula}, matched against the imaginary-root
multiplicities of Theorem~\ref{thm:im-root-gen}) gives the
denominator function
\begin{equation}\label{eq:WKB-match}
\Phi(z) \;=\; \sum_{w\in W^{(2)}(\HB)} \det(w)\biggl[e^{-2\pi i (w\rho,z)}
- \sum_{a\in \HB\cap \R_{>0}\cP_{II}} m(a) e^{-2\pi i(w(\rho+a),z)}\biggr]
\end{equation}
\[
= e^{-2\pi i(\rho, z)}\prod_{\alpha\in\Delta_{+}}
(1 - e^{-2\pi i(\alpha, z)})^{\mathrm{mult}(\alpha)}
\;=\; \tfrac{1}{64}\Delta_{5}(2 Z)
\]
per Lorgat 2020 Theorem 3, valid for $z = z_{3}f_{2}+z_{2}f_{3}+z_{1}
f_{-2}\in\Omega(\HB)$ and $Z = \begin{pmatrix}z_{1}&z_{2}\\ z_{2}&z_{3}
\end{pmatrix}\in H_{2}$ the Siegel half-space.
\end{theorem}

\begin{proof}[Chain-level proof]
Combine (i) Proposition~\ref{prop:shadow-formula} (the shadow is the
chiral-CE complex of $\g \otimes \Omega^{0,\ast}(E)$ tensored with the
$K3$-fibre zero-mode algebra); (ii) Corollary~\ref{cor:HA-from-geometry}
(the lattice is $\HB = \Lambda^{(1,1)}\oplus[2]\simeq
\Z f_{2}\oplus \Z f_{3}\oplus \Z f_{-2}$); (iii) Theorem~\ref{thm:im-root-gen}
plus Proposition~\ref{prop:tau-9} (the imaginary-root multiplicities
are the Fourier coefficients of $\phi_{0,1}$, with $\tau(a_{0}) = 9$);
(iv) the Weyl vector $\rho = \tfrac{1}{2}(\delta_{1}+\delta_{2}+\delta_{3})
= f_{2} - \tfrac{1}{2}f_{3} + f_{-2}$ as in Lorgat 2020 \S4 and Remark 4.1
of the source paper. The Weyl group $W^{(2)}(\HB)$ acts on $\HB\otimes\R$ by
reflections in $\delta_{1},\delta_{2},\delta_{3}$ (Gram matrix
$\mathrm{diag}(2)-2\cdot(\mathbf{1}-I)$), and by Lorgat 2020 Lemma 3 this
is precisely the reflection-anti-invariance action of $\Delta_{5}$. The
character formula of the GKM theory (Borcherds 1988; Lorgat 2020 Theorem
3) equates the Weyl-Kac sum with the infinite product
$e^{-2\pi i(\rho,z)}\prod_{\alpha\in\Delta_{+}}(1-e^{-2\pi i(\alpha,z)}
)^{\mathrm{mult}(\alpha)}$; this is a theorem once the root-system
structure is in place.

The matching with $\tfrac{1}{64}\Delta_{5}(2Z)$ uses Lorgat 2020 \S4
Lemma 4 (Fourier expansion of $\tfrac{1}{64}\Delta_{5}(2Z)$ in terms of
$W^{(2)}(\HB)$-sum with $m(a), \tau(a)$-coefficients matching those of
our imaginary-root multiplicities). The final step (identifying the
Weyl-Kac sum with the Fourier expansion) is an analytic fact about
Siegel modular forms on $\mathrm{Sp}_{4}(\Z)$, holding chain-level up
to the Borcherds-lift transgression which we flag below as conjectural
at the level of full Siegel modularity.
\end{proof}

\begin{remark}[What is conjectural]\label{rem:conjectural-scope}
\ClaimStatusConjectured\
The Siegel modularity of $\tfrac{1}{64}\Delta_{5}(2Z)$ on
$\mathrm{Sp}_{4}(\Z)$, i.e.\ the full transformation law under
$\mathrm{Sp}_{4}(\Z)$ with the multiplier system
$\nu_{\Delta_{5}}:\mathrm{Sp}_{4}(\Z)\to \C^{\times}$ found by Maass,
is \emph{not} derived in this note. The numerical content (weight $5$,
multiplier system values on the three generator types of
$\mathrm{Sp}_{4}(\Z)$, Fourier expansion) matches as shown above, and
the Heegner--lattice structure forces the invariance under the
Jacobi-parabolic subgroup $\Gamma_{\infty}$; however, the full
Sp-equivariance under the non-parabolic generator $S$ requires the
Borcherds lift transgression from weak Jacobi form to Siegel modular
form, which we have not done at Feynman-coefficient level. Full
Siegel modularity is thus a conjecture at the
$\SpF_{K3,E}\circ \PhiFA_{3}$-level; it is a theorem at the
Lorgat-2020 / Gritsenko-Nikulin-1995 / Borcherds-1995 analytic level.
A future Wave may close this gap by constructing an explicit
transgression on the Costello--Gwilliam factorisation-homology
pushforward side.
\end{remark}

%==============================================================
\section{Synthesis: the full chain}
%==============================================================

\begin{theorem}[$6$d hCS on $K3\times E\leadsto \Delta_{5}$ and $\Gdelta$]
\label{thm:hCS-to-Delta5}
\ClaimStatusDerivation\
Let $X = K3\times E$ with holomorphic volume form $\Omega_{X}=
\Omega_{K3}\wedge\omega_{E}$. The following chain, composed of one
classical datum $S_{\mathrm{cl}}$, one renormalisation $I[L]$, and two
functorial specialisations, produces at its terminus the denominator
of the GKM Lie superalgebra $\Gdelta$ of Lorgat 2020:

\begin{enumerate}
\item[(S)] \textbf{Stage 0 (classical BV).}
$\cA\in\Omega^{0,\ast}(X,\g)$ with
$S_{\mathrm{cl}}[\cA] = \int_{X}\Omega_{X}\wedge(\tfrac12\langle\cA,
\dbar\cA\rangle + \tfrac16\langle\cA,[\cA,\cA]\rangle)$. CME holds
(Lemma~\ref{lem:CME-K3E}).

\item[(Q)] \textbf{Stage 1 (renormalisation to hFA).} Costello RG flow
at scale $L$ with propagator $P^{K3}_{\varepsilon, L}\cdot P^{E}_{\varepsilon,L}$;
output is the $E_{3}$-holomorphic factorisation algebra
$\PhiFA_{3}(X) = \Obs^{q}_{\hCS_{6}}(X)$
(Definition~\ref{def:PhiFA-K3E}).

\item[(K)] \textbf{Stage 2A ($K3$-fibre pushforward).} Factorisation
homology along $K3$:
$\int_{K3}\PhiFA_{3}(X)|_{E}$ gives an $\Eone$-chiral algebra on $E$
with Beilinson--Drinfeld chiral-CE structure tensored with the
$K3$-fibre Wilson zero-mode $\g\oplus\g[-2]$
(Proposition~\ref{prop:shadow-formula}).

\item[(L)] \textbf{Stage 2B (lattice identification).}
The total cohomology lattice
$\Lambda^{(5,21)} = \Muk\otimes H^{\ast}(E)$ has a canonical
Heegner-signature-$(3,2)$ slice $\LamF$
(Proposition~\ref{prop:heegner-slice}), with hyperbolic sublattice
$\HB\simeq\Lambda^{(1,1)}\oplus[2]$
(Corollary~\ref{cor:HA-from-geometry}).
$W^{(2)}(\HB)$-reflection group acts on $\HB\otimes\R$.

\item[(I)] \textbf{Stage 2C (imaginary-root multiplicity).}
The $K3$-fibre partition function at $\g=\C$ is the K3-elliptic-genus
weak Jacobi form $\phi_{0,1}=Z_{K3}/2$; its Fourier coefficients
$f(n, l)$ are the multiplicities $\mathrm{mult}(\alpha)$ of positive
roots $\alpha = (n,l,m)$ (Theorem~\ref{thm:im-root-gen}). The cusp
$\tau(a_{0})=9$ identity is the
$r^{1/2}$-Jacobi-triple-product-coefficient of $\phi_{5,1}$
(Proposition~\ref{prop:tau-9}).

\item[(D)] \textbf{Stage 3 (denominator identification).}
The Weyl-Kac-Borcherds character formula on $\Gdelta$
(real even simple roots $\{\delta_{1},\delta_{2},\delta_{3}\}$
with Gram matrix $\mathrm{diag}(2)-2(\mathbf{1}-I)$; imaginary simple
roots from Stage 2C; Weyl group $W^{(2)}(\HB)$) evaluates the
denominator at the trivial $1$-dim module, producing
$\tfrac{1}{64}\Delta_{5}(2Z)$ (Theorem~\ref{thm:denominator-match}).
\end{enumerate}

Stage 0 through Stage 2C are chain-level derivations valid on the
Costello--Gwilliam hCS$_{6}$-avatar. Stage 3 (denominator identification)
is chain-level valid up to the Borcherds-lift transgression that
upgrades the $\Gamma_{\infty}$-invariance to full $\mathrm{Sp}_{4}(\Z)
$-equivariance; the latter is flagged \ClaimStatusConjectured\ at the
hCS$_{6}$-derivation level (Remark~\ref{rem:conjectural-scope}).
\end{theorem}

\begin{conjecture}[Siegel modularity of $\Delta_{5}$ from hCS$_{6}$]
\label{conj:siegel-modularity}
The Siegel modularity of $\tfrac{1}{64}\Delta_{5}(2Z)$ under
$\mathrm{Sp}_{4}(\Z)$ with multiplier $\nu_{\Delta_{5}}$ is a
\emph{consequence} of the hCS$_{6}$-derivation above,
conditional on the Borcherds-lift transgression identity that
upgrades $\Gamma_{\infty}$-invariance of the Stage 2C factorisation-homology
pushforward to full $\mathrm{Sp}_{4}(\Z)$-equivariance. Concretely:
Stage 2B (Proposition~\ref{prop:heegner-slice}) constructs the
$\mathrm{SO}^{+}(\LamF)$-action on the Heegner-slice lattice; the
conjectural content is that the factorisation-homology pushforward
$\SpF_{K3, E}\circ \PhiFA_{3}$ intertwines the
$\mathrm{Sp}_{4}(\Z)$-action via $\wedge^{2}$ (Lorgat 2020 Lemma 1) with
the Weyl-action of $O(\HB)_{+}$ on the chiral-algebra on $E$, forcing
Siegel modularity of the denominator function with the Maass multiplier
$\nu_{\Delta_{5}}$.
\end{conjecture}

%==============================================================
\section{Abelian-Heisenberg first-order verification}
%==============================================================

\begin{proposition}[Abelian verification at $\g = \fu(1)^{22}$]
\label{prop:abelian-check}
\ClaimStatusTheorem\
For $\g = \fu(1)^{\oplus 22}$ (rank-$22$ abelian: K3 Mukai latttice
rank) the hCS$_{6}$-theory on $K3\times E$ is a free
Gaussian theory. The partition function at this abelian level is
\[
Z^{\mathrm{abel}}_{\hCS_{6}}(K3\times E; \tau) \;=\;
\bigl(\det{}'(\dbar_{K3\times E})\bigr)^{-11} \;=\; \theta_{\Muk}(\tau)
\cdot [\text{Jacobi}(\tau, z)]
\]
where $\det'$ is the regularised zero-mode-removed determinant, and
$\theta_{\Muk}(\tau)$ is the Mukai-lattice theta series on $E_{\tau}$:
\[
\theta_{\Muk}(\tau) \;=\; \sum_{v\in \Muk} \exp(\pi i\tau \cdot
\langle v, v\rangle_{\Muk}).
\]
The lattice theta of the Mukai lattice (signature $(4, 20)$) is the
leading-order chain-level manifestation of the $K3$-partition function.
\end{proposition}

\begin{proof}[Proof sketch]
At abelian $\g = \fu(1)^{22}$, the BV action is purely Gaussian:
$S = \tfrac12\int \Omega \wedge A_{0,1}\dbar A_{0,1}$. Path integral
reduces to Gaussian-determinant times zero-mode sum. The zero-modes of
$\dbar$ on $\Omega^{0,1}(K3,\R^{22})$ are parametrised by
$H^{0,1}(K3)\otimes \R^{22} = 0$ (since $H^{0,1}(K3)=0$ for K3), while
$H^{0,0}\otimes \R^{22}$ and $H^{0,2}\otimes \R^{22}$ contribute
$22+22 = 44$ constant modes summing to a lattice sum over the Mukai
lattice extended by $E$-cohomology. The result is the Mukai-lattice
theta series. For the full K3 $\times$ E partition function with
imaginary-root correction one passes to $\g = \fsl_{2}$ non-abelian
hCS$_{6}$, where the cubic vertex generates $K3$-fibre Wilson
one-loop corrections (Cycle 5), producing the $\phi_{0,1}$
multiplicity factor.
\end{proof}

\begin{remark}[Lattice theta with imaginary-root correction]
\label{rem:lattice-theta-correction}
The Lorgat 2020 product expansion
$\tfrac{1}{64}\Delta_{5}(Z) = e^{\pi i(z_{1}+z_{2}+z_{3})}\prod_{
(n,l,m)>0}(1-\exp(2\pi i(nz_{1}+lz_{2}+mz_{3})))^{f(nm,l)}$
(Lorgat 2020 \S6 Theorem 4) is the full imaginary-root-corrected
lattice theta of $\HB$ with signature $(2,1)$. The exponents $f(nm,l)$
are the Fourier coefficients of $\phi_{0,1}=Z_{K3}/2$; the abelian
Mukai-lattice theta of Proposition~\ref{prop:abelian-check} is its
leading-order $\hbar^{0}$ contribution, and the
$\phi_{0,1}$-exponents encode the ``$K3$-imaginary-root correction''
generated at one loop by $\g = \fsl_{2}$.
\end{remark}

%==============================================================
\section{Conclusion and cache appendages}
%==============================================================

\begin{remark}[Summary of the five attack-heal cycles]\label{rem:summary}
\begin{enumerate}
\item[C1.] The ambient theory is Witten hCS on CY$_{3}$, not Costello's 4d
twist. BV field $\cA\in \Omega^{0,\ast}(K3\times E,\g)$; CME holds.
\item[C2.] Observables form an $E_{3}$-hFA on $K3\times E$, not an
$E_{1}$-chiral algebra on the ambient; the $E_{1}$-chiral shadow is
Stage-$2$.
\item[C3.] The $K3$-fibre pushforward $\SpF_{K3, E}$ is Lurie--Francis
factorisation homology; on $K3$ (with $\chi(\cO_{K3})=2, c_{2}(K3)=24
\neq 0$) the wedge-vanishing argument of $\C^{3}$ fails, and the
$K3$-fibre generates a non-trivial chiral shadow on $E$.
\item[C4.] The lattice $\LamF\simeq \Lambda^{(1,1)}\oplus \Lambda^{(1,1)}
\oplus [2]$ of Lorgat 2020 \S3 is the Heegner-signature-$(3,2)$ slice
of the total cohomology $\Lambda^{(5,21)} = \Muk\otimes
H^{\ast}(E,\Z)$-type lattice; $\HB$ is its hyperbolic sublattice.
\item[C5.] The imaginary-root multiplicities of $\Gdelta$ are Fourier
coefficients of $\phi_{0,1}=Z_{K3}/2$; the cusp $\tau(a_{0}) = 9$
identity is the $\eta^{9}$-factor of $\phi_{5, 1/2}=\eta(z_{1})^{9}
\nu_{11}(z_{1},z_{2})$, with $9$ decomposing as $\chi(\cO_{K3})\cdot
\dim\fsl_{2}+\text{triple-product shift}$.
\end{enumerate}
\end{remark}

\begin{remark}[Cache append candidates]
\label{rem:cache-appendage}
Three AP-CY cache candidates for
\texttt{notes/antipatterns\_catalogue.md}:

\emph{(A)} ``hCS$_{6}$ on $K3\times E$ is a chiral algebra'' is a
scope violation: the native observables form an $E_{3}$-hFA, not an
$E_{1}$-chiral algebra. The chiral shadow lives on $E$ via $\SpF_{K3,E}$,
not on the ambient.

\emph{(B)} ``Wedge-square propagator vanishes on CY$_{3}$'' is a
$\C^{3}$-specific statement. On $K3\times E$, the ambient wedge
$\Omega_{X}\wedge\Omega_{X}=0$ still holds, but after $K3$-fibre
factorisation homology the residual $E$-propagator contributions are
non-zero because $\chi(\cO_{K3}) = 2\neq 0$ and $c_{2}(K3) = 24\neq 0$.

\emph{(C)} $\Gdelta$ is the $E_{1}$-chiral shadow's Lie-envelope; its
real simple roots come from the $\HB$-lattice (Heegner slice of $\Muk
\otimes H^{\ast}(E)$); its imaginary simple roots come from the
$K3$-elliptic-genus multiplicities. $\tau(a_{0}) = 9$ is Siegel-
consistent with $\chi(\cO_{K3}) = 2, \chi(\cO_{E})=0$ and the $\eta^{9}$
chain-level origin.
\end{remark}

\begin{remark}[Summary $\leq 400$ words]\label{rem:400-word}
On $X = K3\times E$ with $\Omega_{X} = \Omega_{K3}\wedge\omega_{E}$,
$6$d holomorphic Chern--Simons has BV superfield $\cA\in
\Omega^{0,\ast}(X,\g)$ and classical action $S_{\mathrm{cl}}[\cA] =
\int_{X}\Omega_{X}\wedge(\tfrac12\langle\cA,\dbar\cA\rangle + \tfrac16
\langle\cA,[\cA,\cA]\rangle)$. The classical master equation holds.
Costello renormalisation at scale $L$ with propagator $P^{K3}\cdot
P^{E}$ produces the $E_{3}$-holomorphic factorisation algebra
$\PhiFA_{3}(X) = \Obs^{q}_{\hCS_{6}}(X)$, per Costello--Gwilliam Vol 2
Thm 10.0.1 and CFG 2026 Thm 4.6. At stalks, $U_{E_{3}}(\g[2])$.

Factorisation homology along $K3$ specialises this to an
$E_{1}$-chiral algebra on $E$. The wedge-vanishing
argument making hCS$_{6}$ trivial on $\C^{3}$ fails on $K3$
($\chi(\cO_{K3}) = 2$, $c_{2}(K3)=24$): the $K3$-fibre Wilson insertion
produces a non-trivial $\g[-2]$-summand, and the chiral shadow on $E$
is the Beilinson--Drinfeld chiral-CE complex of $\g\otimes
\Omega^{0,\ast}_{c}(E)$ tensored with the $K3$-fibre zero-mode data.

The signature-$(3,2)$ lattice $\LamF$ of Lorgat 2020 \S3 arises as
the Heegner-slice of $H^{\ast}(K3 \times E, \Z) = \Muk \otimes
H^{\ast}(E, \Z) \simeq \Lambda^{(5,21)}$; the hyperbolic sublattice
$\HB = \Lambda^{(1,1)}\oplus[2]$ is the relevant $(2,1)$-sub for the
GKM structure. The $\wedge^{2}: \mathrm{Sp}_{4}(\Z)/\pm \to
\mathrm{SO}^{+}(\LamF)/\pm$ isomorphism makes Siegel modularity
natural.

The imaginary-root multiplicities $m(a)$, $\tau(a)$ come from the
$K3$-fibre partition function: the weak Jacobi form $\phi_{0,1} =
Z_{K3}/2$ encodes
$\mathrm{mult}(\alpha = (n,l,m))=f(nm,l)$. The cusp identity
$\tau(a_{0}) = 9$ is the $\eta^{9}$-factor of $\phi_{5,1/2}$, which
chain-level decomposes as $\chi(\cO_{K3}) \cdot \dim\fsl_{2} + 3$.
The Weyl--Kac--Borcherds character formula applied to the trivial
module of the GKM Lie superalgebra $\Gdelta$ on root datum
$\{\delta_{1},\delta_{2},\delta_{3}\}\cup\{\phi_{0,1}
\text{-imaginary roots}\}$ produces the denominator function
\[
\Phi(z) = e^{-2\pi i(\rho,z)}\prod_{\alpha\in\Delta_{+}}
(1-e^{-2\pi i(\alpha,z)})^{\mathrm{mult}\alpha} = \tfrac{1}{64}
\Delta_{5}(2Z)
\]
matching Lorgat 2020 Theorem 3. Full $\mathrm{Sp}_{4}(\Z)$-Siegel
modularity under the Maass multiplier is conjectural at the
hCS$_{6}$-pushforward level; it is a theorem at the Borcherds-lift
analytic level.

Scope: the five-cycle chain is chain-level rigorous for Stages 0--2C;
Stage 3 denominator matching is chain-level up to the Borcherds-lift
transgression. The construction produces one of several sibling BKMs
from $\PhiFA_{3}(K3\times E)$; other cycles $\Sigma_{2}$ produce other
denominators (Corollary 3.4 of \texttt{wave12\_u1\_two\_stage\_functor.tex}),
consistent with the CLAUDE.md programme's claim that ``six routes to
$G(K3\times E)$ are six different constructions, not six $\Phi$
applications''.
\end{remark}

%==============================================================
\section*{References (inline)}
%==============================================================
\begin{thebibliography}{99}
\bibitem{Witten1992CS} E.\ Witten, ``Chern--Simons gauge theory as a
string theory'', \emph{Prog.\ Math.}\ 133 (1995), 637--678,
arXiv:hep-th/9207094.
\bibitem{Costello2011} K.\ Costello, \emph{Renormalization and Effective
Field Theory}, AMS Math.\ Surveys 170 (2011).
\bibitem{CGVol1} K.\ Costello, O.\ Gwilliam, \emph{Factorization Algebras
in Quantum Field Theory}, Vol 1, Cambridge (2017).
\bibitem{CGVol2} K.\ Costello, O.\ Gwilliam, \emph{Factorization Algebras
in Quantum Field Theory}, Vol 2, Cambridge (2021).
\bibitem{CostelloFrancisGwilliam2026} K.\ Costello, J.\ Francis, O.\
Gwilliam, ``The $E_{n}$-algebra of observables of $n$-dim holomorphic
CS'', in preparation (2026); structural Theorem 4.6 referenced above.
\bibitem{CostelloLi} K.\ Costello, S.\ Li, ``Quantization of open-closed
BCOV theory I'', arXiv:1505.06703.
\bibitem{Borcherds1995} R.\ Borcherds, ``Automorphic forms with
singularities on Grassmannians'', arXiv:alg-geom/9609022.
\bibitem{GritsenkoNikulin1995} V.\ Gritsenko, V.\ Nikulin, ``Siegel
automorphic form corrections of some Lorentzian Kac-Moody Lie
algebras'', arXiv:alg-geom/9504006.
\bibitem{Gritsenko1999} V.\ Gritsenko, ``Complex vector bundles and
Jacobi forms'', arXiv:math/9906191.
\bibitem{Lorgat2020} R.\ Lorgat, ``A Borcherds lift of the weak Jacobi
form $\phi_{0,1}$, generalized Borcherds--Kac--Moody superalgebras and
the Igusa cusp form $\Delta_{5}$'', unpublished note, April 2020
(this programme's user's own manuscript; see
\texttt{notes/wave13\_source\_paper.txt} for full text).
\bibitem{OberdieckPixton} G.\ Oberdieck, A.\ Pixton, ``Holomorphic
anomaly equations and the Igusa cusp form conjecture'', arXiv 2018.
\bibitem{MukaiLattice} S.\ Mukai, ``On the moduli space of bundles on
K3 surfaces, I'', in \emph{Vector Bundles on Algebraic Varieties},
Bombay 1984.
\bibitem{NikulinK3} V.\ Nikulin, ``Integral symmetric bilinear forms and
some of their applications'', \emph{Math.\ USSR Izv.}\ 14 (1980),
103--167.
\bibitem{BeilinsonDrinfeld} A.\ Beilinson, V.\ Drinfeld,
\emph{Chiral Algebras}, AMS Colloquium Publications 51 (2004).
\bibitem{LurieHA} J.\ Lurie, \emph{Higher Algebra}, author's website.
\bibitem{GriffithsHarris} P.\ Griffiths, J.\ Harris, \emph{Principles of
Algebraic Geometry}, Wiley (1978).
\end{thebibliography}

\end{document}
